\section{Zadanie 1}
Udowodnij, że liczba elementów skończonej algebry (czyli też \(\sigma\)-algebry) jest potęgą dwójki.

\subsection*{Rozwiązanie}
(Tak jak w zadaniu 2) Rozważamy atomy w danej \(\sigma\)-algebrze. Niech n - liczba atomów, zgodnie z warunkami jakie musi spełniać \(\sigma\)-algebra suma każdego podzbioru atomów musi należeć do tej algebry. Liczba podzbiorów zbioru \(n\)-elementowego to \(2^{n}\) zatem liczba elementów skończonej \(\sigma\)-algebry musi być potęgą dwójki CKD
